% ============================================================
%  Basado en la plantilla de la Practica 8 (mismos colores/estilos)
% ============================================================
\documentclass[11pt,a4paper]{article}

% ---------- Idioma y codificación ----------
\usepackage[utf8]{inputenc}
\usepackage[T1]{fontenc}
\usepackage[spanish,es-noshorthands]{babel}

% ---------- Tipografía y layout ----------
\usepackage{lmodern}
\usepackage[sfdefault,light]{roboto}
\usepackage{geometry}
\geometry{a4paper, top=3.1cm, bottom=2.8cm, left=2.4cm, right=2.4cm, headheight=1.6cm}
\usepackage{ragged2e}
\usepackage{enumitem}
\usepackage{graphicx}
\usepackage{float}
\usepackage{booktabs}
\usepackage{array}
\usepackage[most]{tcolorbox}
\usepackage{tikz}
\usetikzlibrary{shadows,calc}
\usepackage{titlesec}
\usepackage{fancyhdr}
\usepackage{fontawesome5}
\usepackage[colorlinks=true,breaklinks=true]{hyperref}
\usepackage{caption}
\usepackage{microtype}
\usepackage{xurl}
\usepackage{listings}
\usepackage{amsmath,amssymb}
\usepackage{needspace}
\usepackage{icomma}

% ---------- Paleta de colores (tema "coolnight" de tu terminal Kitty) ----------
\definecolor{cNavy}{HTML}{011423}
\definecolor{cTealDeep}{HTML}{064B5A}
\definecolor{cTeal}{HTML}{097289}
\definecolor{cBlue}{HTML}{0FC5ED}
\definecolor{cCyan}{HTML}{0FC5ED}
\definecolor{cMint}{HTML}{44FFB1}
\definecolor{cPurple}{HTML}{A277FF}
\definecolor{cRed}{HTML}{E52E2E}
\definecolor{cBg}{HTML}{EAF6F8}
\definecolor{cText}{HTML}{01111E}
\definecolor{cOutBg}{HTML}{F4F4F4}

\hypersetup{
  linkcolor=cTealDeep,
  citecolor=cTealDeep,
  urlcolor=cBlue,
  pdftitle={Distribuciones exponencial y normal: lista de ejercicios resuelta},
  pdfauthor={Miguel Deza}
}

\color{cText}
\emergencystretch=2.5em

% ---------- Encabezado / pie de página ----------
\pagestyle{fancy}
\fancyhf{}
\renewcommand{\headrulewidth}{0pt}
\renewcommand{\footrulewidth}{0pt}
\fancyhead[L]{\small\color{cTealDeep}\textbf{\faGraduationCap\ Maestría en Ciencia de Datos}}
\fancyhead[R]{\small\color{cTeal}Distribuciones Exponencial y Normal}
\fancyfoot[C]{%
  \begin{tikzpicture}[remember picture]
    \node[circle, fill=cTealDeep, minimum size=1.7em, inner sep=0pt] (pg) {\color{white}\small\thepage};
  \end{tikzpicture}%
}
\renewcommand{\headrule}{{\color{cCyan}\hrule height 1.4pt}\vspace{1pt}{\color{cBlue}\hrule height 0.5pt}}

% ---------- Estilo de secciones ----------
\titleformat{\section}
  {\normalfont\Large\bfseries\color{cTealDeep}}
  {\tikz[baseline=-0.6ex]{\node[fill=cCyan, text=white, rounded corners=1.5pt, inner sep=3pt, font=\bfseries]{\thesection};}}
  {0.6em}{}
  [{\vspace{-0.3em}\color{cCyan}\titlerule[1.2pt]\vspace{0.3em}}]

\titleformat{\subsection}
  {\normalfont\large\bfseries\color{cTeal}}
  {\color{cTeal}\thesubsection}{0.6em}{}

\titlespacing*{\section}{0pt}{1.4em}{0.8em}
\titlespacing*{\subsection}{0pt}{1.1em}{0.5em}

% ---------- Figura sin marco ----------
\newcommand{\figuraMarcada}[3]{%
  \begin{figure}[H]
    \centering
    \includegraphics[width=#2\linewidth]{#1}
    \caption{#3}
  \end{figure}
}

\captionsetup{labelfont={bf,color=cTealDeep}, font=small, textfont=it}

% ---------- Bloques de código R y salida de consola ----------
\lstdefinestyle{styleR}{
  language=R,
  backgroundcolor=\color{cBg},
  basicstyle=\ttfamily\small\color{cText},
  keywordstyle=\color{cTealDeep}\bfseries,
  commentstyle=\color{gray}\itshape,
  stringstyle=\color{cTeal},
  numbers=none,
  breaklines=true,
  frame=l,
  framerule=2.2pt,
  rulecolor=\color{cCyan},
  xleftmargin=8pt,
  framexleftmargin=6pt,
  aboveskip=6pt, belowskip=6pt,
  showstringspaces=false
}

\lstdefinestyle{styleOut}{
  backgroundcolor=\color{cOutBg},
  basicstyle=\ttfamily\scriptsize\color{cText},
  breaklines=true,
  frame=l,
  framerule=2.2pt,
  rulecolor=\color{cTealDeep},
  xleftmargin=8pt,
  framexleftmargin=6pt,
  aboveskip=2pt, belowskip=6pt,
  showstringspaces=false
}

\newcommand{\codigoR}[1]{\lstinputlisting[style=styleR]{#1}}

% ---------- Ejercicios ----------
\usepackage{enumitem}
\setlist{itemsep=1pt, topsep=3pt}
\graphicspath{{img/}}
\setlength{\parskip}{3pt}
\newcommand{\R}{\mathbb{R}}

\newtcolorbox{enunciado}[1][]{enhanced, breakable, colback=cBg, colframe=cTealDeep,
  boxrule=0pt, leftrule=3pt, arc=2pt, left=8pt, right=8pt, top=5pt, bottom=5pt,
  fonttitle=\bfseries, coltitle=white, colbacktitle=cTealDeep,
  title={Ejercicio #1}}
\newtcolorbox{resultado}{enhanced, breakable, colback=cMint!12, colframe=cTealDeep,
  boxrule=0.5pt, arc=2pt, left=8pt, right=8pt, top=4pt, bottom=4pt}

\newif\ifinicioseccion
\newcommand{\ejercicio}[2]{%
  \ifinicioseccion\global\inicioseccionfalse\else\clearpage\fi
  \phantomsection
  \addcontentsline{toc}{subsection}{Ejercicio #1}%
  \begin{enunciado}[#1]#2\end{enunciado}\noindent}
% portada de sección: página completa con fondo oscuro
\newif\ifanexo
\newcommand{\seccion}[4]{%
  \clearpage\newgeometry{margin=0cm}
  \refstepcounter{section}\phantomsection
  \addcontentsline{toc}{section}{\protect\numberline{\thesection}\protect\faIcon{#1}\ #2}%
  \thispagestyle{empty}\noindent
  \begin{tikzpicture}[remember picture, overlay]
    \fill[cNavy] (current page.south west) rectangle (current page.north east);
    \fill[cTealDeep, opacity=0.55] (current page.north west) ++(1.5,-1) circle (4.2);
    \fill[cCyan, opacity=0.14] (current page.north west) ++(4.2,-0.4) circle (2.2);
    \fill[cPurple, opacity=0.26] (current page.south east) ++(-1.2,1.4) circle (3.6);
    \fill[cMint, opacity=0.13] (current page.south east) ++(-3.6,0.3) circle (2.0);
    \draw[cCyan, opacity=0.25, line width=0.8pt] (current page.center) ++(5.2,4.6) circle (1.1);
    \node[anchor=center] at ($(current page.center)+(0,5.2)$)
      {\color{cTealDeep!80!cCyan}\bfseries\fontsize{150}{150}\selectfont \thesection};
    \node[circle, fill=cTealDeep, draw=cCyan, line width=1.3pt, minimum size=2.4cm] (ic)
      at ($(current page.center)+(0,1.6)$) {};
    \node at (ic) {\color{cMint}\fontsize{30}{30}\selectfont\faIcon{#1}};
    \node[anchor=center] at ($(current page.center)+(0,-0.4)$)
      {\color{cCyan}\bfseries\fontsize{11}{13}\selectfont \ifanexo A N E X O\else S E C C I Ó N\fi};
    \node[anchor=north, text width=14cm, align=center] (tit) at ($(current page.center)+(0,-0.9)$)
      {\hyphenpenalty=10000 {\color{white}\bfseries\fontsize{30}{36}\selectfont #3}};
    \draw[cCyan, line width=1.4pt] ($(tit.south)+(-3.2,-0.3)$) -- ++(6.4,0);
    \node[anchor=north] at ($(tit.south)+(0,-0.6)$)
      {\color{cMint}\fontsize{13}{16}\selectfont #4};
  \end{tikzpicture}%
  \clearpage\restoregeometry}

% figura a la izquierda y su código R a la derecha
% (las líneas de cada bloque se leen de rangos.tex, generado desde practica-10-distribuciones-exponencial-normal.R)
\expandafter\def\csname r@ejE1-servidor\endcsname{35-44}
\expandafter\def\csname r@ejE2-api\endcsname{47-58}
\expandafter\def\csname r@ejE3-sesiones\endcsname{61-73}
\expandafter\def\csname r@ejE4-iot\endcsname{76-88}
\expandafter\def\csname r@ejE5-call-center\endcsname{91-99}
\expandafter\def\csname r@ejE6-clics\endcsname{102-110}
\expandafter\def\csname r@ejN1-error-ml\endcsname{113-123}
\expandafter\def\csname r@ejN2-salario\endcsname{126-134}
\expandafter\def\csname r@ejN3-sensor\endcsname{137-148}
\expandafter\def\csname r@ejN4-logit\endcsname{151-164}
\expandafter\def\csname r@ejN5-clustering\endcsname{167-179}
\expandafter\def\csname r@ejN6-batchnorm\endcsname{182-193}
\expandafter\def\csname r@ejN7-lidar\endcsname{196-206}
\expandafter\def\csname r@ejN8-regresion\endcsname{209-222}
\expandafter\def\csname r@ejN9-recomendacion\endcsname{225-240}

\newcommand{\figcod}[2]{%
  \par\medskip\noindent
  \begin{minipage}[c]{0.47\linewidth}\centering
    \includegraphics[width=\linewidth]{#1}\captionof{figure}{#2}
  \end{minipage}\hfill
  \begin{minipage}[c]{0.5\linewidth}
    \edef\figopts{linerange={\csname r@#1\endcsname}}%
    \expandafter\lstinputlisting\expandafter[\figopts, style=styleR, basicstyle=\ttfamily\scriptsize\color{cText}]{practica-10-distribuciones-exponencial-normal.R}
  \end{minipage}\par\medskip}
\newcommand{\res}[1]{\begin{resultado}\textbf{\color{cTealDeep}Resultado.} #1\end{resultado}}

% ============================================================
\begin{document}

% -------------------- PORTADA --------------------
\begin{titlepage}
\newgeometry{margin=0cm}
\begin{tikzpicture}[remember picture, overlay]
  \fill[cNavy] (current page.south west) rectangle (current page.north east);
  \fill[cPurple, opacity=0.28] (current page.north east) ++(-1,1) circle (3.2);
  \fill[cCyan, opacity=0.22] (current page.north east) ++(-2.6,0.4) circle (2.1);
  \fill[cTealDeep, opacity=0.55] (current page.south west) ++(1,-1) circle (4.0);
  \fill[cMint, opacity=0.16] (current page.south west) ++(2.4,-0.2) circle (2.3);

  \node[circle, fill=cTealDeep, draw=cCyan, line width=1.3pt, minimum size=2.5cm] (logo) at ($(current page.north)+(0,-2.9)$) {};
  \node at (logo) {\color{cMint}\fontsize{28}{28}\selectfont\faIcon{chart-area}};

  \node[anchor=north, text width=14cm, align=center] at ($(current page.north)+(0,-4.7)$) {
    {\color{cCyan}\bfseries\fontsize{12}{14}\selectfont \faDatabase\ \ MAESTRÍA EN CIENCIA DE DATOS}
  };

  \node[anchor=center, text width=15cm, align=center] at ($(current page.center)+(0,1.2)$) {
    {\color{white}\bfseries\fontsize{32}{40}\selectfont Distribuciones\\ Exponencial y Normal}\\[0.35em]
    {\color{cMint}\bfseries\fontsize{20}{24}\selectfont \faIcon{chart-area}\ en Ciencia de Datos}
  };

  \draw[cCyan, line width=1.4pt] ($(current page.center)+(-3.6,-0.9)$) -- ++(7.2,0);

  \node[anchor=center, text width=13cm, align=center] at ($(current page.center)+(0,-1.8)$) {
    {\color{white!85!black}\fontsize{11}{14}\selectfont
     Solución de la lista de ejercicios:\\
     tiempos entre eventos, cuantiles, errores de modelos,\\ estandarización y truncamiento}
  };

  \node[anchor=south, text width=14cm, align=center] at ($(current page.south)+(0,2.6)$) {
    \begin{tabular}{r l}
      \color{cCyan}\bfseries \faUser\ Elaborado por: & \color{white} Miguel Deza \\[3pt]
      \color{cCyan}\bfseries \faGraduationCap\ Programa: & \color{white} Maestría en Ciencia de Datos -- UNSA \\[3pt]
      \color{cCyan}\bfseries \faBook\ Unidad: & \color{white} I -- Fundamentos Matemáticos \\[3pt]
      \color{cCyan}\bfseries \faChalkboardTeacher\ Docente: & \color{white} Carlos Antonio Abanto Valle \\[3pt]
      \color{cCyan}\bfseries \faIcon{calendar-alt}\ Fecha: & \color{white} 25 de octubre de 2026 \\
    \end{tabular}
  };
\end{tikzpicture}
\end{titlepage}
\restoregeometry

\setcounter{tocdepth}{2}
\tableofcontents
\newpage

% ============================================================
\seccion{hourglass-half}{Distribución exponencial}{Distribución exponencial}{Ejercicios E1 -- E6}
% ============================================================

\ejercicio{E1}{\textbf{Tiempo entre fallas de un servidor.} En un sistema cloud, el tiempo entre fallas
críticas sigue una distribución exponencial con media de 8 horas. Sea $T$ el tiempo (en horas) hasta la
próxima falla. (a) Determine el parámetro $\lambda$. (b) Calcule la probabilidad de que el sistema falle
en menos de 5 horas. (c) Calcule la probabilidad de que \emph{no falle} en las próximas 12 horas.
(d) Determine la mediana del sistema.}

\[
f(t)=\lambda e^{-\lambda t},\qquad F(t)=1-e^{-\lambda t},\qquad P(T>t)=e^{-\lambda t},\qquad t\ge0
\]
(a)
\[
E[T]=\frac1\lambda=8 \quad\Rightarrow\quad \lambda=\frac18=0,125\ \text{fallas/hora}
\]
(b)
\[
P(T<5)=1-e^{-5/8}=1-0,535261=0,4647
\]
(c)
\[
P(T>12)=e^{-12/8}=e^{-1,5}=0,2231
\]
(d)
\[
F(m)=\frac12 \;\Rightarrow\; e^{-m/8}=\frac12 \;\Rightarrow\; m=8\ln2=5,545\ \text{h}
\]

\figcod{ejE1-servidor}{Densidad de $T\sim\text{Exp}(1/8)$.}

\res{$\lambda=0,125$; \ $P(T<5)=0,4647$; \ $P(T>12)=0,2231$; \ mediana $=5,545$ h.}

\ejercicio{E2}{\textbf{Tiempo de respuesta de una API.} Una API tiene tiempos de respuesta $T$ (en
segundos) con distribución exponencial de parámetro $\lambda=3$. (a) Calcule $P(T<0,1)$.
(b) Calcule $P(T>0,5)$. (c) Determine un valor $t$ tal que $P(T>t)=0,01$.}

(a)
\[
P(T<0,1)=1-e^{-3(0,1)}=1-e^{-0,3}=1-0,740818=0,2592
\]
(b)
\[
P(T>0,5)=e^{-3(0,5)}=e^{-1,5}=0,2231
\]
(c)
\[
e^{-3t}=0,01 \;\Rightarrow\; -3t=\ln0,01 \;\Rightarrow\; t=\frac{\ln100}{3}=\frac{4,605170}{3}=1,535\ \text{s}
\]

\figcod{ejE2-api}{Densidad de $T\sim\text{Exp}(3)$.}

\res{$P(T<0,1)=0,2592$; \ $P(T>0,5)=0,2231$; \ $t=1,535$ s.}

\ejercicio{E3}{\textbf{Duración de sesiones de usuario.} El tiempo (en minutos) que un usuario permanece
en una plataforma sigue $T\sim\text{Exp}(\lambda=0,25)$. (a) Probabilidad de que una sesión dure más
de 10 minutos. (b) Probabilidad de que dure entre 2 y 6 minutos. (c) Calcule $E[T]$ y
$\operatorname{Var}(T)$. (d) Determine el cuantil 90\,\% del tiempo de sesión.}

(a)
\[
P(T>10)=e^{-0,25(10)}=e^{-2,5}=0,0821
\]
(b)
\[
P(2<T<6)=e^{-0,25(2)}-e^{-0,25(6)}=e^{-0,5}-e^{-1,5}=0,606531-0,223130=0,3834
\]
(c)
\[
E[T]=\frac1\lambda=\frac{1}{0,25}=4\ \text{min},\qquad \operatorname{Var}(T)=\frac{1}{\lambda^2}=\frac{1}{0,0625}=16\ \text{min}^2
\]
(d)
\[
1-e^{-0,25\,q}=0,90 \;\Rightarrow\; e^{-0,25\,q}=0,10 \;\Rightarrow\; q_{0,90}=\frac{\ln10}{0,25}=9,210\ \text{min}
\]

\figcod{ejE3-sesiones}{Densidad (arriba) y función de distribución (abajo) de $T\sim\text{Exp}(0,25)$.}

\res{$P(T>10)=0,0821$; \ $P(2<T<6)=0,3834$; \ $E[T]=4$, \ $\operatorname{Var}(T)=16$; \ $q_{0,90}=9,21$ min.}

\ejercicio{E4}{\textbf{Detección de anomalías en IoT.} El tiempo entre eventos anómalos sigue una
distribución exponencial con media 30 minutos. (a) Probabilidad de que transcurran dos horas sin
anomalías. (b) Probabilidad de detectar una anomalía antes de 15 minutos. (c) Determine $t$ tal que la
probabilidad de \emph{no anomalía} hasta $t$ minutos sea 0,20.}

\[
\lambda=\frac{1}{30}\ \text{por minuto}
\]
(a)
\[
P(T>120)=e^{-120/30}=e^{-4}=0,0183
\]
(b)
\[
P(T<15)=1-e^{-15/30}=1-e^{-0,5}=0,3935
\]
(c)
\[
P(T>t)=e^{-t/30}=0,20 \;\Rightarrow\; t=-30\ln0,20=30\ln5=48,28\ \text{min}
\]

\figcod{ejE4-iot}{Densidad (arriba) y supervivencia (abajo) de $T\sim\text{Exp}(1/30)$.}

\res{$P(T>120)=0,0183$; \ $P(T<15)=0,3935$; \ $t=48,28$ min.}

\ejercicio{E5}{\textbf{Modelo de abandono en call center.} El tiempo hasta que un usuario cuelga si nadie
atiende es $T\sim\text{Exp}(\lambda=0,4)$ (en minutos). (a) Probabilidad de que abandone antes de 2
minutos. (b) Probabilidad de que siga esperando más de 5 minutos. (c) Determine el percentil 75.}

(a)
\[
P(T<2)=1-e^{-0,4(2)}=1-e^{-0,8}=1-0,449329=0,5507
\]
(b)
\[
P(T>5)=e^{-0,4(5)}=e^{-2}=0,1353
\]
(c)
\[
1-e^{-0,4\,q}=0,75 \;\Rightarrow\; q_{0,75}=\frac{\ln4}{0,4}=\frac{1,386294}{0,4}=3,466\ \text{min}
\]

\figcod{ejE5-call-center}{Densidad de $T\sim\text{Exp}(0,4)$.}

\res{$P(T<2)=0,5507$; \ $P(T>5)=0,1353$; \ $P_{75}=3,466$ min.}

\ejercicio{E6}{\textbf{Tiempo entre clics en una campaña.} Una campaña genera clics siguiendo un proceso de
Poisson con tasa $\lambda=12$ clics por hora. El tiempo entre clics es exponencial (en minutos).
(a) Probabilidad de que entre dos clics haya menos de 2 minutos. (b) Probabilidad de que pasen más de
10 minutos sin clic. (c) Determine la mediana del tiempo entre clics.}

\[
\lambda=\frac{12}{60}=0,2\ \text{clics/min},\qquad T\sim\text{Exp}(0,2)
\]
(a)
\[
P(T<2)=1-e^{-0,4}=1-0,670320=0,3297
\]
(b)
\[
P(T>10)=e^{-2}=0,1353
\]
(c)
\[
e^{-0,2\,m}=\frac12 \;\Rightarrow\; m=\frac{\ln2}{0,2}=3,466\ \text{min}
\]

\figcod{ejE6-clics}{Densidad del tiempo entre clics.}

\res{$P(T<2)=0,3297$; \ $P(T>10)=0,1353$; \ mediana $=3,466$ min.}

% ============================================================
\seccion{chart-area}{Distribución normal}{Distribución normal}{Ejercicios N1 -- N9}
% ============================================================

\ejercicio{N1}{\textbf{Error de predicción en un modelo de ML.} El error de un modelo de regresión sigue
$E\sim N(0,2^2)$. (a) Probabilidad de un error mayor a 3. (b) Probabilidad de que el error esté entre
$-1$ y 1. (c) Calcule el valor esperado de $|E|$.}

\[
Z=\frac{E-0}{2}\sim N(0,1)
\]
(a)
\[
P(E>3)=P\!\left(Z>\tfrac32\right)=1-\Phi(1,5)=1-0,933193=0,0668
\]
(b)
\[
P(-1<E<1)=P(-0,5<Z<0,5)=2\Phi(0,5)-1=2(0,691462)-1=0,3829
\]
(c)
\[
E|E|=2\int_0^{\infty}e\,\frac{1}{\sigma\sqrt{2\pi}}e^{-e^2/(2\sigma^2)}\,de
=\frac{2}{\sigma\sqrt{2\pi}}\Big[-\sigma^2e^{-e^2/(2\sigma^2)}\Big]_0^{\infty}
=\frac{2\sigma}{\sqrt{2\pi}}=\sigma\sqrt{\frac2\pi}
\]
\[
E|E|=2\sqrt{\frac2\pi}=2(0,797885)=1,596
\]

\figcod{ejN1-error-ml}{Densidad de $E$ (arriba) y de $|E|$ (abajo).}

\res{$P(E>3)=0,0668$; \ $P(-1<E<1)=0,3829$; \ $E|E|=1,596$.}

\ejercicio{N2}{\textbf{Imputación de datos faltantes.} El salario mensual (en miles) sigue
$X\sim N(5;\ 1,2^2)$. (a) Probabilidad de que un salario sea mayor a 7. (b) Probabilidad de que esté
entre 4 y 6. (c) Determine el cuantil 95\,\% de $X$.}

(a)
\[
P(X>7)=P\!\left(Z>\frac{7-5}{1,2}\right)=P(Z>1,667)=1-0,952210=0,0478
\]
(b)
\[
P(4<X<6)=P\!\left(\frac{4-5}{1,2}<Z<\frac{6-5}{1,2}\right)=P(-0,833<Z<0,833)=2(0,797672)-1=0,5953
\]
(c)
\[
x_{0,95}=\mu+z_{0,95}\,\sigma=5+1,6449(1,2)=6,974
\]

\figcod{ejN2-salario}{Densidad de $X\sim N(5;\ 1,2^2)$.}

\res{$P(X>7)=0,0478$; \ $P(4<X<6)=0,5953$; \ $x_{0,95}=6,974$ mil.}

\ejercicio{N3}{\textbf{Detección de outliers en sensores.} Un sensor mide temperatura con
$T\sim N(50,4^2)$. (a) Probabilidad de una lectura mayor a 60. (b) Probabilidad de una lectura menor
a 40. (c) Intervalo central al 90\,\%. (d) Umbral del percentil 99\,\%.}

(a)
\[
P(T>60)=P\!\left(Z>\frac{60-50}{4}\right)=P(Z>2,5)=1-0,993790=0,0062
\]
(b)
\[
P(T<40)=P(Z<-2,5)=1-\Phi(2,5)=0,0062
\]
(c)
\[
50\pm z_{0,95}(4)=50\pm1,6449(4)=50\pm6,579 \;\Rightarrow\; [43,42;\ 56,58]
\]
(d)
\[
t_{0,99}=50+z_{0,99}(4)=50+2,3263(4)=59,31
\]

\figcod{ejN3-sensor}{Colas (arriba) e intervalo central al 90\,\% (abajo).}

\res{$P(T>60)=P(T<40)=0,0062$; \ intervalo 90\,\%: $[43,42;\ 56,58]$; \ $P_{99}=59,31$.}

\ejercicio{N4}{\textbf{Puntajes de un clasificador probabilístico.} El logit sigue
$Z\sim N(1,5;\ 0,7^2)$ y $p=\dfrac{1}{1+e^{-Z}}$. (a) Probabilidad de que el puntaje sea negativo.
(b) Probabilidad de que sea mayor que 2. (c) Use una aproximación de primer orden para estimar $E[p]$.}

(a)
\[
P(Z<0)=P\!\left(U<\frac{0-1,5}{0,7}\right)=\Phi(-2,143)=0,0161,\qquad U\sim N(0,1)
\]
(b)
\[
P(Z>2)=P\!\left(U>\frac{2-1,5}{0,7}\right)=1-\Phi(0,714)=1-0,762475=0,2375
\]
(c) $\sigma(z)=\dfrac{1}{1+e^{-z}}$, \ Taylor alrededor de $\mu=1,5$:
\[
p=\sigma(Z)\approx\sigma(\mu)+\sigma'(\mu)(Z-\mu)
\quad\Rightarrow\quad E[p]\approx\sigma(\mu)+\sigma'(\mu)\underbrace{E[Z-\mu]}_{0}=\sigma(1,5)
\]
\[
E[p]\approx\frac{1}{1+e^{-1,5}}=\frac{1}{1+0,223130}=0,8176
\]
\textbf{Comparación:}
\[
\sigma''(\mu)=\sigma(1-\sigma)(1-2\sigma)=-0,094730
\]
\[
E[p]\approx0,8176+\tfrac12(-0,094730)(0,49)=0,7944\ \ (\text{2.\textsuperscript{o} orden}),\qquad E[p]_{\text{exacto}}=0,7965
\]

\figcod{ejN4-logit}{Densidad del logit $Z$ (arriba) y de $p=\sigma(Z)$ (abajo).}

\res{$P(Z<0)=0,0161$; \ $P(Z>2)=0,2375$; \ $E[p]\approx\sigma(1,5)=0,8176$ (exacto 0,7965).}

\ejercicio{N5}{\textbf{Resultado de un algoritmo de clustering.} El clúster A tiene $X_A\sim N(20,3^2)$ y
el clúster B tiene $X_B\sim N(27,2^2)$. (a) Probabilidad de que un usuario de A tenga valor mayor que
25. (b) Probabilidad de que un usuario de B tenga valor menor que 25. (c) Si se observa un usuario con
valor 23, ¿a qué clúster pertenece con mayor verosimilitud?}

(a)
\[
P(X_A>25)=P\!\left(Z>\frac{25-20}{3}\right)=P(Z>1,667)=0,0478
\]
(b)
\[
P(X_B<25)=P\!\left(Z<\frac{25-27}{2}\right)=\Phi(-1)=0,1587
\]
(c)
\[
f_A(23)=\frac{1}{3\sqrt{2\pi}}e^{-\frac{(23-20)^2}{2(9)}}=\frac{\varphi(1)}{3}=\frac{0,241971}{3}=0,0807
\]
\[
f_B(23)=\frac{1}{2\sqrt{2\pi}}e^{-\frac{(23-27)^2}{2(4)}}=\frac{\varphi(2)}{2}=\frac{0,053991}{2}=0,0270
\]
\[
\frac{f_A(23)}{f_B(23)}=2,99>1 \quad\Rightarrow\quad A
\]

\figcod{ejN5-clustering}{Densidades de los clústeres A y B.}

\res{$P(X_A>25)=0,0478$; \ $P(X_B<25)=0,1587$; \ el usuario con valor 23 pertenece al clúster A.}

\ejercicio{N6}{\textbf{Normal estandarizada en ML (Batch Normalization).} Sea $X\sim N(\mu,\sigma^2)$ con
$\mu=45$ y $\sigma=8$, y $Z=\dfrac{X-\mu}{\sigma}$. (a) Calcule $P(Z>2)$. (b) Calcule $P(-1<Z<1)$.
(c) Determine el valor original $x$ del cuantil 97,5\,\%.}

(a)
\[
P(Z>2)=1-\Phi(2)=1-0,977250=0,0228
\]
(b)
\[
P(-1<Z<1)=2\Phi(1)-1=2(0,841345)-1=0,6827
\]
(c)
\[
z_{0,975}=1,96 \;\Rightarrow\; x=\mu+z\,\sigma=45+1,96(8)=60,68
\]

\figcod{ejN6-batchnorm}{Normal estándar (arriba) y $X\sim N(45,8^2)$ (abajo).}

\res{$P(Z>2)=0,0228$; \ $P(-1<Z<1)=0,6827$; \ $x_{0,975}=60,68$.}

\ejercicio{N7}{\textbf{Ruido en señales de un auto autónomo.} El ruido en un sensor LIDAR es
$R\sim N(0;\ 1,5^2)$. (a) Probabilidad de que el error supere 2 cm en valor absoluto.
(b) Probabilidad de que esté entre $-1$ y 1 cm. (c) Calcule la desviación media $E[|R|]$.}

(a)
\[
P(|R|>2)=2P\!\left(Z>\frac{2}{1,5}\right)=2\big(1-\Phi(1,333)\big)=2(0,091211)=0,1824
\]
(b)
\[
P(-1<R<1)=2\Phi\!\left(\frac{1}{1,5}\right)-1=2(0,747507)-1=0,4950
\]
(c)
\[
E|R|=\sigma\sqrt{\frac2\pi}=1,5(0,797885)=1,197\ \text{cm}
\]

\figcod{ejN7-lidar}{Densidad de $R$ (arriba) y de $|R|$ (abajo).}

\res{$P(|R|>2)=0,1824$; \ $P(-1<R<1)=0,4950$; \ $E|R|=1,197$ cm.}

\ejercicio{N8}{\textbf{Error de regresión múltiple.} Sea $\varepsilon\sim N(0,5^2)$. (a) Calcule
$P(|\varepsilon|>7)$. (b) Calcule $P(-3<\varepsilon<3)$. (c) Determine los cuartiles $Q_1$ y $Q_3$.}

(a)
\[
P(|\varepsilon|>7)=2\big(1-\Phi(1,4)\big)=2(1-0,919243)=0,1615
\]
(b)
\[
P(-3<\varepsilon<3)=2\Phi(0,6)-1=2(0,725747)-1=0,4515
\]
(c)
\[
z_{0,75}=0,6745 \;\Rightarrow\; Q_1=0-0,6745(5)=-3,372,\qquad Q_3=0+0,6745(5)=3,372
\]
\[
\text{RIC}=Q_3-Q_1=6,745
\]

\figcod{ejN8-regresion}{Probabilidades (arriba) y cuartiles (abajo) de $\varepsilon\sim N(0,5^2)$.}

\res{$P(|\varepsilon|>7)=0,1615$; \ $P(-3<\varepsilon<3)=0,4515$; \ $Q_1=-3,372$, \ $Q_3=3,372$.}

\ejercicio{N9}{\textbf{Puntuaciones en un modelo de recomendación.} Las calificaciones predichas siguen
$Y\sim N(3,8;\ 0,6^2)$. (a) Calcule $P(Y>5)$. (b) Calcule $P(3\le Y\le4)$. (c) Determine el percentil
99\,\%. (d) Discuta cualitativamente el efecto de truncar $Y$ al intervalo $[1,5]$.}

(a)
\[
P(Y>5)=P\!\left(Z>\frac{5-3,8}{0,6}\right)=P(Z>2)=0,0228
\]
(b)
\[
P(3\le Y\le4)=\Phi\!\left(\frac{4-3,8}{0,6}\right)-\Phi\!\left(\frac{3-3,8}{0,6}\right)
\]
\[
=\Phi(0,333)-\Phi(-1,333)=0,630559-0,091211=0,5393
\]
(c)
\[
y_{0,99}=3,8+2,3263(0,6)=5,196>5
\]
(d) \textbf{Truncamiento en $[1,5]$:}
\[
P(1\le Y\le5)=\Phi(2)-\Phi(-4,667)=0,977248,\qquad
f_T(y)=\frac{f(y)}{0,977248},\quad 1\le y\le5
\]
\[
E[Y\mid1\le Y\le5]=3,8+0,6\,\frac{\varphi(-4,667)-\varphi(2)}{0,977248}=3,767,\qquad
\sigma_T=0,565<0,6
\]
Se elimina el 2,28\,\% de masa sobre 5 (bajo 1 es $\approx10^{-6}$) y se reparte en $[1,5]$: la media
baja de 3,80 a 3,77, la desviación de 0,600 a 0,565 y el $P_{99}$ pasa de 5,20 (fuera de escala) a 4,91.

\figcod{ejN9-recomendacion}{Probabilidades (arriba) y densidad original frente a truncada (abajo).}

\res{$P(Y>5)=0,0228$; \ $P(3\le Y\le4)=0,5393$; \ $P_{99}=5,196$; \ truncada: media 3,767, desv. 0,565.}

\end{document}
